\documentclass[11pt,a4paper]{article}

% ── Packages ──────────────────────────────────────────────────────────────────
\usepackage[a4paper, margin=2.5cm]{geometry}
\usepackage{fontenc}
\usepackage{inputenc}
\usepackage{microtype}
\usepackage{parskip}
\usepackage{titlesec}
\usepackage{titletoc}
\usepackage{tocloft}
\usepackage{enumitem}
\usepackage{booktabs}
\usepackage{longtable}
\usepackage{tabularx}
\usepackage{array}
\usepackage{multirow}
\usepackage{xcolor}
\usepackage{mdframed}
\usepackage{listings}
\usepackage{fancyhdr}
\usepackage{graphicx}
\usepackage{hyperref}
\usepackage{amsmath}
\usepackage{amssymb}
\usepackage{float}
\usepackage{caption}
\usepackage{subcaption}
\usepackage{appendix}
\usepackage{tcolorbox}
\tcbuselibrary{breakable, skins}


% ── Colors ────────────────────────────────────────────────────────────────────
\definecolor{criticalred}{RGB}{180,40,40}
\definecolor{criticallight}{RGB}{253,235,235}
\definecolor{importantyellow}{RGB}{160,110,0}
\definecolor{importantlight}{RGB}{255,248,220}
\definecolor{minorgray}{RGB}{80,80,100}
\definecolor{minorlight}{RGB}{240,240,248}
\definecolor{fixblue}{RGB}{20,80,160}
\definecolor{fixlight}{RGB}{232,242,255}
\definecolor{codecolor}{RGB}{248,248,252}
\definecolor{coderule}{RGB}{200,200,220}
\definecolor{sectionblue}{RGB}{15,60,130}
\definecolor{linkcolor}{RGB}{20,80,160}
\definecolor{notegreen}{RGB}{30,110,60}
\definecolor{notelight}{RGB}{234,248,238}

% ── Hyperref setup ────────────────────────────────────────────────────────────
\hypersetup{
  colorlinks=true,
  linkcolor=linkcolor,
  urlcolor=linkcolor,
  citecolor=linkcolor,
  pdftitle={Real-Time Tennis Oddsmaking and Integrity Monitoring},
  pdfauthor={Team},
  pdfsubject={Big Data System Architecture and Implementation Plan},
  bookmarksnumbered=true,
  bookmarksopen=true
}

% ── Section heading style ─────────────────────────────────────────────────────
\titleformat{\section}{\large\bfseries\color{sectionblue}}{{\thesection}}{1em}{}[\vspace{-4pt}\rule{\linewidth}{0.4pt}]
\titleformat{\subsection}{\normalsize\bfseries\color{sectionblue}}{\thesubsection}{1em}{}
\titleformat{\subsubsection}{\normalsize\bfseries}{\thesubsubsection}{1em}{}

% ── Header / footer ───────────────────────────────────────────────────────────
\pagestyle{fancy}
\fancyhf{}
\fancyhead[L]{\small\textit{Real-Time Tennis Oddsmaking \& Integrity Monitoring}}
\fancyhead[R]{\small\textit{Big Data Architecture}}
\fancyfoot[C]{\thepage}
\renewcommand{\headrulewidth}{0.4pt}

% ── Code listing style ────────────────────────────────────────────────────────
\lstset{
  basicstyle=\ttfamily\small,
  backgroundcolor=\color{codecolor},
  frame=single,
  framerule=0.5pt,
  rulecolor=\color{coderule},
  breaklines=true,
  breakatwhitespace=true,
  showstringspaces=false,
  tabsize=2,
  captionpos=b,
  xleftmargin=6pt,
  xrightmargin=6pt,
  aboveskip=8pt,
  belowskip=8pt
}

% ── Custom boxes ──────────────────────────────────────────────────────────────
\newtcolorbox{criticalbox}[1][]{
  breakable, enhanced,
  colback=criticallight, colframe=criticalred,
  fonttitle=\bfseries\small\color{white},
  title={\textbf{CRITICAL GAP}\ #1},
  coltitle=white, attach boxed title to top left={yshift=-2mm,xshift=4mm},
  boxed title style={colback=criticalred, colframe=criticalred},
  top=6pt, left=8pt, right=8pt, bottom=6pt
}

\newtcolorbox{importantbox}[1][]{
  breakable, enhanced,
  colback=importantlight, colframe=importantyellow,
  fonttitle=\bfseries\small,
  title={\textbf{IMPORTANT GAP}\ #1},
  coltitle=importantyellow, attach boxed title to top left={yshift=-2mm,xshift=4mm},
  boxed title style={colback=importantlight, colframe=importantyellow},
  top=6pt, left=8pt, right=8pt, bottom=6pt
}

\newtcolorbox{minorbox}[1][]{
  breakable, enhanced,
  colback=minorlight, colframe=minorgray,
  fonttitle=\bfseries\small,
  title={\textbf{MINOR GAP}\ #1},
  coltitle=minorgray, attach boxed title to top left={yshift=-2mm,xshift=4mm},
  boxed title style={colback=minorlight, colframe=minorgray},
  top=6pt, left=8pt, right=8pt, bottom=6pt
}

\newtcolorbox{fixbox}{
  breakable, enhanced,
  colback=fixlight, colframe=fixblue,
  left=8pt, right=8pt, top=5pt, bottom=5pt,
  fontupper=\small
}

\newtcolorbox{notebox}{
  breakable, enhanced,
  colback=notelight, colframe=notegreen,
  left=8pt, right=8pt, top=5pt, bottom=5pt,
  fontupper=\small
}

% ── Table column types ────────────────────────────────────────────────────────
\newcolumntype{L}[1]{>{\raggedright\arraybackslash}p{#1}}
\newcolumntype{C}[1]{>{\centering\arraybackslash}p{#1}}
\newcolumntype{R}[1]{>{\raggedleft\arraybackslash}p{#1}}

% ── Misc ──────────────────────────────────────────────────────────────────────
\setlength{\parskip}{6pt}
\setlength{\parindent}{0pt}
\captionsetup{font=small, labelfont=bf}

% ─────────────────────────────────────────────────────────────────────────────
\begin{document}

% ── Title page ────────────────────────────────────────────────────────────────
\begin{titlepage}
  \centering
  \vspace*{2cm}
  {\huge\bfseries\color{sectionblue} Real-Time Tennis Oddsmaking\\[6pt] and Integrity Monitoring\par}
  \vspace{1cm}
  \rule{\linewidth}{1pt}
  \vspace{4pt}
  {\large A Big Data System Architecture and Implementation Plan\\
  \textit{Revised Edition --- All Architectural Gaps Addressed}\par}
  \vspace{4pt}
  \rule{\linewidth}{1pt}
  \vspace{1.5cm}
  {\normalsize Team Members: Add Names and NetIDs\par}
  \vspace{0.5cm}
  {\normalsize April 28, 2026\par}
  \vspace{2cm}
  \begin{tcolorbox}[colback=fixlight, colframe=fixblue, width=0.85\linewidth]
    \centering\small
    \textbf{Revision Note.}
    This document incorporates all 24 architectural gaps identified in the
    original proposal review. New or modified content is present throughout
    every section. Gap annotations (\textcolor{criticalred}{\textbf{CRITICAL}},
    \textcolor{importantyellow}{\textbf{IMPORTANT}},
    \textcolor{minorgray}{\textbf{MINOR}}) mark each
    addition so reviewers can locate changes quickly.
  \end{tcolorbox}
\end{titlepage}

% ── Table of contents ─────────────────────────────────────────────────────────
\tableofcontents
\newpage

% =============================================================================
\section{Purpose of This Document}
% =============================================================================

This document defines the final architecture, implementation plan, data contracts, pipeline design,
modeling approach, testing strategy, and presentation plan for the big data project titled
\textit{Real-Time Tennis Oddsmaking and Integrity Monitoring}. The goal is to provide a single
reference that guides development, team coordination, cloud deployment, report writing, and final
demonstration.

The project follows a practical Lambda-style architecture deployed using two Amazon EC2 instances
and an Amazon S3 data lake. EC2 Instance~1 acts as the ingestion, orchestration, streaming, batch
processing, and machine learning node. EC2 Instance~2 acts as the serving and application node.
Amazon~S3 is the durable data lake storing raw data, cleaned data, feature tables, model
artifacts, baseline profiles, benchmark outputs, and final report assets.

% =============================================================================
\section{Project Thesis and Big Data Justification}
% =============================================================================

\subsection{Project Thesis}

The project builds a real-time analytics system for professional tennis supporting two related use
cases from the same underlying event stream. First, it computes live win probabilities for
simulated tennis matches using point-by-point streaming events. Second, it monitors match-integrity
risk by comparing live player behavior against historical baselines and flagging statistically
unusual deviations.

The central thesis is that live sports prediction and integrity monitoring can be unified into one
data engineering system because both require granular event streams, historical baselines, rolling
features, low-latency scoring, and serving layers that expose results to end users.

\subsection{Why This Is a Data Engineering Problem}

This project is primarily a data engineering problem because the system must move, transform,
store, and serve data across multiple stages. The core work is not only building a machine
learning model, but building the infrastructure around the model. The system must ingest raw
historical files, prepare replayable event streams, process streaming data, build feature tables,
train models, store outputs, and serve results through APIs and dashboards.

Core data engineering challenges include:
\begin{itemize}[nosep]
  \item maintaining a clean data lake layout in S3;
  \item converting historical point-level data into live replay streams;
  \item processing Kafka events with Spark Structured Streaming;
  \item generating batch features and player baselines;
  \item coordinating offline and online pipelines via Airflow DAG dependencies;
  \item persisting serving-ready outputs in PostgreSQL and optional Redis;
  \item exposing results through FastAPI and React.
\end{itemize}

\subsection{Why This Is a Big Data Project Rather Than a Small Local Script}

A small script could train a static model on a subset of tennis matches but would not demonstrate
the distributed and system-level requirements of this project. This project involves high-volume
historical processing, synthetic stream amplification, real-time event ingestion, stateful
streaming computations, batch retraining, and separate serving responsibilities.

The big data nature comes from:
\begin{itemize}[nosep]
  \item \textbf{Volume:} historical data expanded into a larger synthetic archive;
  \item \textbf{Velocity:} point events replayed as live streams through Kafka;
  \item \textbf{Variety:} match-level, point-level, player-level, model, and risk-event data coexist;
  \item \textbf{Complexity:} rolling features, baselines, anomaly scores, and live odds computed consistently;
  \item \textbf{Deployment:} multiple services across two cloud instances and S3.
\end{itemize}

\subsection{Target Scale Envelope}

\begin{itemize}[nosep]
  \item Replay hundreds to thousands of simulated matches depending on EC2 resources;
  \item process point events through Kafka topics partitioned by match identifier;
  \item store raw and processed data in Parquet format in S3;
  \item maintain low-latency serving of current match states through PostgreSQL and optional Redis;
  \item expose dashboards that refresh frequently enough to feel real-time for demonstration.
\end{itemize}

\subsection{Numerical Benchmark Targets}

\begin{table}[H]
\centering
\begin{tabularx}{\linewidth}{L{5.5cm}X}
\toprule
\textbf{Metric} & \textbf{Target} \\
\midrule
Streaming throughput & At least 1,000 point events per second for MVP; higher as stretch goal \\
End-to-end latency & Under 5 seconds from event production to dashboard visibility \\
Batch processing & Complete feature generation on selected dataset within demo window \\
Model output latency & Streaming inference completed within each Spark micro-batch \\
Serving read latency & API responses under 500\,ms in normal demo conditions \\
Dashboard refresh & 2--5 second refresh interval for live views \\
\bottomrule
\end{tabularx}
\caption{Benchmark targets}
\end{table}

% =============================================================================
\section{Executive Architecture Decision}
% =============================================================================

\subsection{Final Architectural Choice}

The final architecture is a practical Lambda-style architecture using:

\begin{itemize}[nosep]
  \item Amazon S3 as the central data lake;
  \item EC2 Instance~1 as the processing node;
  \item EC2 Instance~2 as the serving and application node;
  \item Apache Kafka for live event ingestion;
  \item Apache Spark for batch processing and streaming analytics;
  \item Apache Airflow for workflow orchestration;
  \item PostgreSQL for serving current match, odds, and risk-event data;
  \item optional Redis for low-latency latest-state caching;
  \item FastAPI for backend APIs;
  \item React for the frontend dashboard.
\end{itemize}

\subsection{Why This Choice Was Made}

This architecture balances ambition, cost control, and implementation feasibility. Fully managed
services such as Amazon MSK, Amazon EMR, and managed OpenSearch would be closer to a production
cloud architecture but are cost-prohibitive for a student project. A purely local system would be
cheaper but would not demonstrate realistic cloud deployment.

The two-EC2 architecture provides a strong middle ground: real services running in AWS, separated
processing from serving, S3 as a real cloud data lake, and manageable operational complexity.

% =============================================================================
\section{System Overview}
% =============================================================================

\subsection{High-Level Components}

\begin{table}[H]
\centering
\begin{tabularx}{\linewidth}{L{3.5cm}X}
\toprule
\textbf{Component} & \textbf{Responsibility} \\
\midrule
S3 Data Lake & Stores raw, cleaned, features, models, baselines, logs, and benchmark outputs \\
EC2 Instance~1 & Runs ingestion, Kafka, Spark, ML workflows, and Airflow \\
EC2 Instance~2 & Runs PostgreSQL, optional Redis, FastAPI, React, and optional Nginx \\
Kafka & Buffers and distributes simulated live point events \\
Spark Batch & Builds curated datasets, features, baselines, and model training inputs \\
Spark Streaming & Processes live point events, computes rolling features, and performs online scoring \\
Airflow & Schedules and monitors batch workflows and model publication jobs \\
PostgreSQL & Stores serving-ready current match state, odds history, and risk events \\
Redis & Optional cache for latest match state and live dashboard reads \\
FastAPI & Provides API endpoints to the dashboard \\
React & Provides the user-facing live odds and risk exploration interface \\
\bottomrule
\end{tabularx}
\caption{High-level system components}
\end{table}

\subsection{Canonical Flow}

\begin{enumerate}[nosep]
  \item Raw tennis datasets are uploaded to S3.
  \item EC2 Instance~1 reads raw data from S3.
  \item Spark batch jobs clean data and build curated datasets.
  \item Feature engineering jobs create offline feature tables and player baselines.
  \item Model training jobs train odds and risk models.
  \item Model artifacts and baselines are published back to S3 via a validated two-phase publish.
  \item A replay producer reads prepared point sequences and sends events to Kafka.
  \item Spark Structured Streaming consumes Kafka events.
  \item Streaming jobs compute rolling features and apply latest models.
  \item Current odds and risk outputs are written to PostgreSQL and optionally Redis on EC2 Instance~2.
  \item FastAPI exposes serving data through API endpoints.
  \item React dashboard displays live odds, match drill-downs, risk events, and benchmark summaries.
\end{enumerate}

% =============================================================================
\section{Scope Definition}
% =============================================================================

\subsection{What Will Be Real}

\begin{itemize}[nosep]
  \item S3 data lake layout with zone manifests;
  \item Kafka-based event replay with explicit topic configuration;
  \item Spark batch jobs for cleaning and feature generation;
  \item Spark streaming job for live processing with dead-letter handling;
  \item baseline odds model;
  \item baseline anomaly/risk scoring model;
  \item PostgreSQL serving schema with indexes and archival;
  \item FastAPI backend with bearer token authentication;
  \item React dashboard;
  \item benchmark scripts and result exports.
\end{itemize}

\subsection{What May Be Simplified}

\begin{itemize}[nosep]
  \item Number of features may be reduced to a practical MVP set;
  \item model choices may use scikit-learn or Spark MLlib depending on complexity;
  \item Redis may be optional if PostgreSQL performance is sufficient;
  \item Airflow initially orchestrates one MVP DAG for the offline pipeline;
  \item streaming scale may be lower than a full production sports-betting workload.
\end{itemize}

\subsection{What Will Not Be Claimed}

\begin{itemize}[nosep]
  \item Real betting odds suitable for financial wagering;
  \item anomaly flags as proof of match-fixing or misconduct;
  \item synthetic replay exactly matching live professional tennis operations;
  \item two-instance deployment at production scale;
  \item legally or commercially validated models.
\end{itemize}

% =============================================================================
\section{Repository Structure}
% =============================================================================

\begin{lstlisting}
tennis-big-data-project/
README.md
.github/
  workflows/
    ci.yml               # linting, unit tests, compose-validate
docker-compose.processing.yml
docker-compose.serving.yml
docker-compose.dev.yml   # single-machine dev with shared bridge network
.env.example
data/
  sample/
contracts/
  point_event_schema.json    # enforced on raw ingestion
  match_state_schema.json
  risk_event_schema.json
infra/
  ec2/
    kafka_setup.sh       # topic creation with explicit partition config
    spark_defaults.conf  # Iceberg catalog + micro-batch tuning
  s3/
    create_layout.sh     # creates prefixes + zone_metadata.json
  iam/
    ec2-processing-role-policy.json
    ec2-serving-role-policy.json
  nginx/
producer/
  replay_producer.py
  config.py
spark/
  batch/
    historical_ingestion.py
    curated_build.py
    feature_build.py
    baseline_build.py
    train_models.py
    publish_model.py     # two-phase validated model publish
  streaming/
    streaming_processor.py
    sinks.py
    scoring.py
    dead_letter.py       # malformed-message handling
airflow/
  dags/
    nightly_batch_pipeline.py   # A->B->C->D->E->publish->F
    replay_demo_pipeline.py
api/
  app/
    main.py
    middleware/
      auth.py            # bearer token check
    routes/
    db/
    schemas/
frontend/
  src/
    pages/
    components/
    services/
sql/
  serving_schema.sql
  indexes.sql            # idx_odds_history_match_ts + others
  archival.sql           # pg_cron archival job
tests/
  unit/
  integration/
  e2e/
reports/
  benchmark_results/
\end{lstlisting}

% =============================================================================
\section{Technology Decisions}
% =============================================================================

\subsection{Core Stack}

\begin{table}[H]
\centering
\begin{tabularx}{\linewidth}{L{4cm}L{4cm}X}
\toprule
\textbf{Layer} & \textbf{Technology} & \textbf{Reason} \\
\midrule
Cloud Storage & Amazon S3 & Durable data lake and artifact store \\
Lakehouse Table Format & Apache Iceberg (Hadoop catalog) & Schema evolution, snapshots, incremental reads on curated/features zones \\
Compute & Two Amazon EC2 instances & Cost-controlled cloud deployment \\
Streaming Ingestion & Apache Kafka & Event buffering and replay simulation \\
Batch Processing & Apache Spark & Scalable feature engineering \\
Streaming Processing & Spark Structured Streaming & Stateful live event processing \\
Orchestration & Apache Airflow & Scheduled workflows with explicit DAG dependencies \\
Serving Database & PostgreSQL & Reliable relational serving store \\
Cache & Redis (optional) & Fast latest-state reads \\
Backend & FastAPI & Lightweight Python API layer \\
Frontend & React & Interactive dashboard \\
Containerization & Docker Compose & Reproducible deployment \\
Secrets Management & AWS SSM Parameter Store & Shared secrets across both EC2 instances \\
\bottomrule
\end{tabularx}
\caption{Core technology stack}
\end{table}

\subsection{Iceberg Catalog Configuration}

\begin{importantbox}[--- Gap \#09 addressed]
Iceberg requires a catalog backend. For this two-EC2 student deployment the \textbf{Hadoop catalog
backed by S3} is used. No external catalog service is required. The following configuration must
be present in \texttt{infra/ec2/spark\_defaults.conf} and referenced in
\texttt{docker-compose.processing.yml}:
\end{importantbox}

\begin{lstlisting}[language=bash, caption={spark-defaults.conf --- Iceberg Hadoop catalog}]
spark.sql.extensions=org.apache.iceberg.spark.extensions.IcebergSparkSessionExtensions
spark.sql.catalog.iceberg=org.apache.iceberg.spark.SparkCatalog
spark.sql.catalog.iceberg.type=hadoop
spark.sql.catalog.iceberg.warehouse=s3://tennis-big-data/iceberg/
spark.sql.defaultCatalog=iceberg
\end{lstlisting}

\begin{notebox}
\textbf{Iceberg scope.} Iceberg is applied only to the \texttt{cleaned/} and \texttt{features/}
zones. The \texttt{raw/} zone remains plain S3 objects. The \texttt{baselines/} and
\texttt{models/} zones use versioned Parquet files, not Iceberg tables, because they are written
once per training run rather than updated incrementally.
\end{notebox}

\subsection{Explicit Technology Rejections}

\begin{itemize}[nosep]
  \item \textbf{Amazon MSK:} rejected due to cost for student version;
  \item \textbf{Amazon EMR:} optional for stretch benchmarking only;
  \item \textbf{Elasticsearch/OpenSearch:} replaced by PostgreSQL risk-event serving;
  \item \textbf{Cassandra:} replaced by PostgreSQL and optional Redis;
  \item \textbf{Kubernetes:} Docker Compose sufficient for two EC2 instances.
\end{itemize}

% =============================================================================
\section{Environment Split}
% =============================================================================

\subsection{Local Development}

\begin{notebox}
\textbf{Gap \#23 addressed.} A dedicated \texttt{docker-compose.dev.yml} runs all services on a
single machine using a shared Docker bridge network. This eliminates configuration drift between
developers and removes the need to manage two sets of private IPs during local testing. Elastic
IPs are applied to both EC2 instances so private IPs remain stable across reboots.
\end{notebox}

Local development responsibilities:
\begin{itemize}[nosep]
  \item writing and testing Spark transformations on samples;
  \item testing producer payloads;
  \item developing FastAPI endpoints;
  \item developing React pages;
  \item running schema validation and unit tests.
\end{itemize}

\subsection{Cloud Deployment}

\textbf{EC2 Instance~1 --- Processing Node.}
Kafka, producer, Spark, Airflow, ML pipelines.
Recommended type: \texttt{m5.xlarge} (4~vCPUs, 16\,GB RAM).

\begin{criticalbox}[--- Gap \#10 addressed]
EC2-1 runs Kafka, Spark, and Airflow concurrently on 16\,GB RAM. To prevent OOM during model
training:
\begin{itemize}[nosep]
  \item Mount a 30\,GB EBS scratch volume at \texttt{/mnt}; set
    \texttt{spark.local.dir=/mnt/spark-temp} in \texttt{spark-defaults.conf}.
  \item Configure 8\,GB swap: \texttt{sudo fallocate -l 8G /swapfile \&\& sudo mkswap
    /swapfile \&\& sudo swapon /swapfile}.
  \item When running Pipeline~E (model training), \textbf{pause the Kafka consumer and replay
    producer} to free RAM. Document this in the Airflow DAG via task dependencies.
\end{itemize}
\end{criticalbox}

\textbf{EC2 Instance~2 --- Serving and Application Node.}
PostgreSQL, optional Redis, FastAPI, React, optional Nginx.
Recommended type: \texttt{t3.large} or \texttt{m5.large}.

\subsection{Networking and Security Group Setup}

Both EC2 instances must be in the same AWS region, VPC, and preferably the same private subnet.
EC2-1 communicates with EC2-2 over private IP addresses only.

\begin{criticalbox}[--- Gap \#05 addressed: IAM role separation]
Two separate IAM roles must be created and documented in \texttt{infra/iam/}:

\textbf{ec2-processing-role} (attached to EC2-1):
\begin{itemize}[nosep]
  \item \texttt{s3:GetObject}, \texttt{s3:PutObject}, \texttt{s3:DeleteObject} on all prefixes;
  \item \texttt{ssm:GetParameter} on \texttt{/tennis/*};
  \item \texttt{ec2:DescribeInstances} for instance self-identification.
\end{itemize}

\textbf{ec2-serving-role} (attached to EC2-2):
\begin{itemize}[nosep]
  \item \texttt{s3:GetObject} on \texttt{models/} and \texttt{baselines/} prefixes \textbf{only};
  \item \texttt{ssm:GetParameter} on \texttt{/tennis/*};
  \item \textbf{No S3 write permissions} --- EC2-2 must never write to the data lake.
\end{itemize}

EC2-2 security group must allow PostgreSQL (5432) and Redis (6379) traffic only from EC2-1's
private security group. Public inbound must be restricted to SSH from developer IPs and
port~80/443 or~8000 for the demo.
\end{criticalbox}

\subsection{Shared Secrets Management}

\begin{importantbox}[--- Gap \#18 addressed]
All shared secrets (\texttt{POSTGRES\_PASSWORD}, \texttt{REDIS\_HOST}, \texttt{API\_KEY}, etc.)
are stored in \textbf{AWS Systems Manager Parameter Store} as \texttt{SecureString} values under
the \texttt{/tennis/} prefix. Both EC2 instances read from SSM at startup:

\begin{lstlisting}[language=bash]
export POSTGRES_PASSWORD=$(aws ssm get-parameter \
  --name /tennis/postgres_password \
  --with-decryption \
  --query Parameter.Value \
  --output text)
\end{lstlisting}

This prevents secrets from appearing in git, \texttt{.env} files synced between instances, or
EC2 instance metadata. The \texttt{.env.example} file references SSM parameter names, not values.
\end{importantbox}

% =============================================================================
\section{Data Contracts}
% =============================================================================

\subsection{Point Event Contract}

\begin{lstlisting}[ caption={Point event Kafka message schema}]
{
  "event_id":       "string",
  "match_id":       "string",
  "event_ts":       "timestamp (UTC ISO-8601)",
  "tournament":     "string",
  "surface":        "string",
  "player_a":       "string",
  "player_b":       "string",
  "server":         "string",
  "receiver":       "string",
  "point_winner":   "string",
  "set_number":     1,
  "game_number":    4,
  "point_number":   23,
  "score_context":  "string",
  "serve_number":   1,
  "rally_length":   6,
  "is_break_point": false,
  "is_double_fault":false,
  "is_ace":         false
}
\end{lstlisting}

\subsection{Current Match State Contract}

\begin{lstlisting}[ caption={Current match state serving contract}]
{
  "match_id":            "string",
  "updated_at":          "timestamp (UTC ISO-8601)",
  "player_a":            "string",
  "player_b":            "string",
  "player_a_win_prob":   0.62,
  "player_b_win_prob":   0.38,
  "momentum_player":     "string",
  "momentum_score":      0.14,
  "set_score":           "string",
  "game_score":          "string",
  "points_played":       87,
  "status":              "live"
}
\end{lstlisting}

\subsection{Risk Event Contract}

\begin{lstlisting}[ caption={Risk event contract}]
{
  "risk_event_id":   "string",
  "match_id":        "string",
  "event_ts":        "timestamp (UTC ISO-8601)",
  "player":          "string",
  "opponent":        "string",
  "risk_score":      0.83,
  "risk_level":      "high",
  "primary_signal":  "serve_win_rate_drop",
  "feature_deviation_summary": {
    "serve_win_rate_delta":    -0.31,
    "rally_length_delta":       4.2,
    "pressure_point_delta":    -0.22
  },
  "explanation": "Current rolling serve performance is far below historical baseline."
}
\end{lstlisting}

\subsection{Schema Governance Rules}

\begin{itemize}[nosep]
  \item Every event must include a unique \texttt{event\_id}.
  \item Every event must include \texttt{match\_id} for partitioning and joins.
  \item Timestamp fields must use UTC ISO-8601.
  \item New fields must be backward compatible where possible.
  \item Null handling rules must be defined for every feature.
  \item Schema changes must be documented and increment the schema version.
\end{itemize}

\begin{importantbox}[--- Gap \#24 addressed: schema enforcement on raw ingestion]
Raw zone files are plain S3 objects (no Iceberg). To prevent mixed-schema Parquet files after
pipeline changes, \texttt{historical\_ingestion.py} validates every incoming file against
\texttt{contracts/point\_event\_schema.json} before writing. Any column additions require a schema
version bump (\texttt{v1}, \texttt{v2}, \ldots) and a full raw-zone re-ingest. Failed validation
routes files to \texttt{s3://tennis-big-data/raw/quarantine/}.
\end{importantbox}

% =============================================================================
\section{Data Lake Design}
% =============================================================================

\subsection{Zones}

\begin{table}[H]
\centering
\begin{tabularx}{\linewidth}{L{3cm}L{3cm}X}
\toprule
\textbf{Zone} & \textbf{Format} & \textbf{Description} \\
\midrule
\texttt{raw/} & CSV / plain S3 & Original source files; schema-validated on ingest \\
\texttt{cleaned/} & Parquet (Iceberg) & Standardized, validated records \\
\texttt{features/} & Parquet (Iceberg) & Offline feature tables \\
\texttt{baselines/} & Parquet (versioned) & Player baseline profiles \\
\texttt{models/} & Pickle/MLlib (versioned) & Trained model artifacts \\
\texttt{stream\_archive/} & Parquet & Archived point events and streaming outputs \\
\texttt{results/} & JSON/CSV & Benchmark and report outputs \\
\bottomrule
\end{tabularx}
\caption{S3 data lake zones}
\end{table}

\subsection{Storage Layout}

\begin{lstlisting}
s3://tennis-big-data/
  raw/
    matches/
    points/
    quarantine/          # schema-validation failures
    zone_metadata.json   # zone manifest
  cleaned/               # Iceberg-managed
    matches/
    points/
    zone_metadata.json
  features/              # Iceberg-managed
    match_features/
    point_features/
    zone_metadata.json
  baselines/
    player_baselines/
    zone_metadata.json
  models/
    odds/
      v1/, v2/, ...
      latest.json        # two-phase publish pointer
    risk/
      v1/, v2/, ...
      latest.json
    zone_metadata.json
  stream_archive/
    point_events/
    scored_events/
    zone_metadata.json
  results/
    benchmarks/
    reports/
    zone_metadata.json
  iceberg/               # Iceberg Hadoop catalog metadata root
\end{lstlisting}

\begin{minorbox}[--- Gap \#19 addressed: zone manifests]
Each zone root contains a \texttt{zone\_metadata.json} file written by
\texttt{infra/s3/create\_layout.sh}. The file records:
\texttt{zone\_name}, \texttt{managed\_by} (\texttt{iceberg}/\texttt{parquet}/\texttt{plain}),
\texttt{owner}, \texttt{created\_at}, \texttt{description}. This prevents confusion when team
members browse the bucket.
\end{minorbox}

\subsection{Partitioning}

\begin{itemize}[nosep]
  \item By year for historical match data;
  \item by tournament or surface for selected feature tables;
  \item by replay date for stream archive data;
  \item by model version for artifacts.
\end{itemize}

\subsection{File Formats}

\begin{itemize}[nosep]
  \item \textbf{CSV:} raw data ingestion only;
  \item \textbf{Parquet:} physical storage for cleaned data, features, stream archives;
  \item \textbf{Apache Iceberg:} table format for curated and feature zones;
  \item \textbf{JSON:} event contracts, zone manifests, model registry pointers;
  \item \textbf{Pickle/MLlib/Joblib:} model artifacts depending on framework.
\end{itemize}

% =============================================================================
\section{Lake vs.\ Serving Responsibilities}
% =============================================================================

\subsection{Data Lake Responsibilities}

S3 is responsible for durable, reproducible storage. It holds the full historical record,
intermediate datasets, offline features, models, baselines, and benchmark evidence.

\subsection{Serving Layer Responsibilities}

The serving layer holds dashboard-ready data. PostgreSQL stores current match states, odds
history, risk events, model versions, and benchmark summaries. Redis may store the latest state
for low-latency reads.

\subsection{Design Rule}

The data lake is the source of truth. The serving layer is optimized for application access. If
serving tables are corrupted or deleted, they must be fully rebuildable from S3 and pipeline
outputs.

% =============================================================================
\section{Serving Data Model}
% =============================================================================

\subsection{PostgreSQL Tables}

\begin{lstlisting}[language=sql, caption={serving\_schema.sql}]
CREATE TABLE matches (
  match_id    TEXT PRIMARY KEY,
  player_a    TEXT NOT NULL,
  player_b    TEXT NOT NULL,
  tournament  TEXT,
  surface     TEXT,
  status      TEXT,
  created_at  TIMESTAMP,
  updated_at  TIMESTAMP
);

CREATE TABLE current_match_state (
  match_id            TEXT PRIMARY KEY REFERENCES matches(match_id),
  updated_at          TIMESTAMP NOT NULL,
  player_a_win_prob   DOUBLE PRECISION,
  player_b_win_prob   DOUBLE PRECISION,
  momentum_player     TEXT,
  momentum_score      DOUBLE PRECISION,
  set_score           TEXT,
  game_score          TEXT,
  points_played       INTEGER
);

CREATE TABLE odds_history (
  id                  BIGSERIAL PRIMARY KEY,
  match_id            TEXT REFERENCES matches(match_id),
  event_ts            TIMESTAMP NOT NULL,
  player_a_win_prob   DOUBLE PRECISION,
  player_b_win_prob   DOUBLE PRECISION,
  momentum_score      DOUBLE PRECISION,
  points_played       INTEGER
);

CREATE TABLE risk_events (
  risk_event_id  TEXT PRIMARY KEY,
  match_id       TEXT REFERENCES matches(match_id),
  event_ts       TIMESTAMP NOT NULL,
  player         TEXT,
  opponent       TEXT,
  risk_score     DOUBLE PRECISION,
  risk_level     TEXT,
  primary_signal TEXT,
  explanation    TEXT
);

CREATE TABLE model_versions (
  model_version  TEXT PRIMARY KEY,
  model_type     TEXT,
  artifact_path  TEXT,
  trained_at     TIMESTAMP,
  metrics_json   JSONB
);
\end{lstlisting}

\begin{criticalbox}[--- Gap \#03 addressed: indexes and archival]
\textbf{Indexes (indexes.sql):}
\begin{lstlisting}[language=sql]
-- Dashboard reads filter by match_id and sort by event_ts DESC
CREATE INDEX idx_odds_history_match_ts
  ON odds_history (match_id, event_ts DESC);

CREATE INDEX idx_risk_events_match_ts
  ON risk_events (match_id, event_ts DESC);

CREATE INDEX idx_risk_events_level
  ON risk_events (risk_level, event_ts DESC);
\end{lstlisting}

\textbf{Archival job (archival.sql --- scheduled via pg\_cron every 15 minutes):}
\begin{lstlisting}[language=sql]
-- Retain only the last 2 hours of odds_history rows
DELETE FROM odds_history
WHERE event_ts < NOW() - INTERVAL '2 hours';

-- Archive risk events older than 24 hours to S3 before deletion
-- (handled by Airflow Pipeline I before DELETE runs)
DELETE FROM risk_events
WHERE event_ts < NOW() - INTERVAL '24 hours';
\end{lstlisting}

\textbf{JDBC batch size} must be set in \texttt{streaming\_processor.py}:
\begin{lstlisting}[language=python]
odds_df.write \
  .format("jdbc") \
  .option("url", POSTGRES_URL) \
  .option("dbtable", "odds_history") \
  .option("batchsize", "500") \
  .mode("append") \
  .save()
\end{lstlisting}
\end{criticalbox}

\subsection{Redis Keys (Optional)}

\begin{lstlisting}
match:{match_id}:latest_state    TTL = 30s
match:{match_id}:latest_odds     TTL = 30s
live_matches                     TTL = 10s
system:latest_benchmark          TTL = 300s
\end{lstlisting}

\begin{importantbox}[--- Gap \#15 addressed: odds\_history size estimate]
For 10 concurrent matches at a 2-second micro-batch trigger over 90 minutes of demo time,
\texttt{odds\_history} accumulates approximately $10 \times 2{,}700 = 27{,}000$ rows. This is
manageable but must be a conscious design decision. The pg\_cron archival job above bounds the
table to the last 2 hours. Live dashboard reads use \texttt{current\_match\_state} (one row per
match) rather than querying \texttt{odds\_history} for the latest value.
\end{importantbox}

\subsection{Serving Retention Strategy}

PostgreSQL retains all demo-period current states, odds histories within the 2-hour rolling
window, and risk events within 24 hours. Redis retains only latest values and is treated as
disposable cache. Long-term archives remain in S3.

% =============================================================================
\section{Synthetic Data Strategy}
% =============================================================================

\subsection{Core Principle}

Synthetic data should amplify scale without inventing unrealistic conclusions. Synthetic replay
preserves the structure of real tennis point sequences while allowing many simulated matches to
run concurrently.

\subsection{Allowed Methods}

\begin{itemize}[nosep]
  \item Replaying historical point sequences with new synthetic match identifiers;
  \item shifting event timestamps to simulate live progression;
  \item sampling matches by tournament, surface, or player;
  \item modest perturbation of non-critical metadata for scale testing;
  \item duplicating replay streams for controlled throughput benchmarks.
\end{itemize}

\subsection{Disallowed Methods}

\begin{itemize}[nosep]
  \item Inventing risk labels and presenting them as real evidence;
  \item claiming synthetic anomalies are confirmed match-fixing;
  \item corrupting tennis scoring rules;
  \item using future points to compute current live features;
  \item producing models that leak target outcomes into live prediction features.
\end{itemize}

\begin{importantbox}[--- Gap \#16 addressed: anomaly training strategy]
The risk model must be trained without invented anomaly labels. The approach is fully
\textbf{unsupervised}:
\begin{enumerate}[nosep]
  \item Compute historical player baselines (mean and standard deviation of each feature
    per player per surface) from the full historical dataset.
  \item During streaming, compute z-scores of current rolling features against the stored
    baseline.
  \item Flag events where $|z| > 2.5$ on two or more key signals simultaneously as
    \texttt{risk\_level = high}.
  \item Optionally train an Isolation Forest on historical feature vectors to score
    global anomalousness alongside the z-score approach.
\end{enumerate}
This approach requires no labeled anomaly examples and is not susceptible to label fabrication.
Document this in \texttt{README.md} under the Modeling section.
\end{importantbox}

\subsection{Documentation Requirement}

Every synthetic generation method must be documented with: input source, transformation rule,
output location, intended purpose, and limitations.

% =============================================================================
\section{Modeling Plan}
% =============================================================================

\subsection{Predictive Odds Model}

The odds model predicts each player's probability of winning given current match state. For the
MVP the model may use Logistic Regression, Random Forest, or Gradient Boosted Trees.

\textbf{Inputs:} current score context, serve win rate, return points won, recent point streak,
surface, player baseline strength, rolling momentum features.

\textbf{Outputs:} player A win probability, player B win probability, model version metadata.

\subsection{Risk / Anomaly Model}

The risk model identifies unusual deviations between current player behavior and historical
baselines using z-score deviation and optional Isolation Forest (see Section~13.2 above).

\textbf{Risk signals:} serve win rate significantly below baseline, double faults above baseline,
unusual drop in pressure-point performance, abnormal rally length pattern, sudden momentum
collapse.

\subsection{Why Spark MLlib Is Not the Sole Path}

Spark MLlib is useful for distributed training but the project may use scikit-learn or another
Python ML library for MVP if the selected dataset is manageable. Spark remains central for
distributed feature generation and streaming analytics.

\subsection{Model Governance Rules}

\begin{itemize}[nosep]
  \item Every model artifact must have a version identifier.
  \item Every model must store its training timestamp.
  \item Evaluation metrics must be saved to \texttt{benchmark\_results/model\_eval.json}.
  \item The streaming job must log which model version was used.
  \item Risk scores must be described as statistical signals, not proof of misconduct.
\end{itemize}

\begin{criticalbox}[--- Gap \#13 addressed: two-phase model publish with validation gate]
\texttt{publish\_model.py} implements a two-phase publish to prevent corrupt or incomplete
model artifacts from reaching the streaming job:
\begin{enumerate}[nosep]
  \item Train and save model to \texttt{s3://tennis-big-data/models/odds/staging/}.
  \item Load the staged artifact and score a fixed test fixture
    (\texttt{data/sample/model\_validation\_fixture.parquet}).
  \item Assert: \texttt{AUC >= 0.55} and \texttt{brier\_score <= 0.30}. If either fails,
    abort --- do not update \texttt{latest.json}.
  \item On success, copy artifact from \texttt{staging/} to \texttt{v\{N\}/} and atomically
    overwrite \texttt{latest.json} with the new version path.
\end{enumerate}
S3 object versioning is enabled on the \texttt{models/} prefix as a recovery fallback.
\end{criticalbox}

\begin{importantbox}[--- Gap \#20 addressed: model evaluation thresholds]
The following minimum thresholds gate model publication. They are deliberately conservative for a
synthetic dataset:
\begin{itemize}[nosep]
  \item Odds model: AUC $\geq 0.55$, Brier score $\leq 0.30$.
  \item Risk model: precision on held-out injected anomalies $\geq 0.50$ (if labeled test set
    available); otherwise, validate that z-score distribution on the test set matches expected
    Gaussian properties.
\end{itemize}
These values are logged to \texttt{results/benchmarks/model\_eval.json} on every training run.
\end{importantbox}

During streaming, the processor loads the latest approved model at startup and refreshes only at
controlled boundaries. Model artifacts are identified by \texttt{latest.json}. A version switch
is logged when it occurs.

% =============================================================================
\section{Feature Design}
% =============================================================================

\subsection{Feature Families}

\begin{table}[H]
\centering
\begin{tabularx}{\linewidth}{L{4cm}X}
\toprule
\textbf{Family} & \textbf{Example Features} \\
\midrule
Serve Dynamics & first serve \%, ace rate, double fault rate, serve points won \\
Return and Rally & return points won, avg.\ rally length, rally length trend \\
Pressure and Clutch & break-point conversion, break-point save rate, deuce performance \\
Momentum and Fatigue & recent points won, streak length, cumulative rally load, momentum index \\
Contextual & surface, tournament, round, player baseline, opponent strength \\
\bottomrule
\end{tabularx}
\caption{Feature families}
\end{table}

\subsection{Rules for Feature Correctness}

\begin{criticalbox}[--- Gap \#07 addressed: point-in-time correctness]
\textbf{Batch features (Pipeline~C):}
\begin{itemize}[nosep]
  \item Sort all events by \texttt{event\_ts} before computing any rolling window.
  \item Apply a \textbf{temporal train/test split}: all matches completed before date~$T$ form
    the training set; matches after $T$ form the test set. Never allow future match data to
    appear in training features.
  \item Use \texttt{Window.rowsBetween(Window.unboundedPreceding, -1)} in Spark window
    functions so the current row's own outcome is never included in its own feature.
  \item Document the split date and window boundaries in \texttt{features/feature\_catalog.md}.
\end{itemize}

\textbf{Streaming features (live):}
\begin{itemize}[nosep]
  \item Streaming features may only use events with \texttt{event\_ts $\leq$ current\_event\_ts}.
  \item Watermark of 30 seconds ensures late events are bounded (see Section~18.3).
  \item Rolling windows use \texttt{withWatermark} before any stateful aggregation.
\end{itemize}
\end{criticalbox}

\begin{itemize}[nosep]
  \item Null values must be explicitly handled (fill with baseline mean or zero as documented).
  \item Feature definitions must be documented and unit-tested.
\end{itemize}

% =============================================================================
\section{Data Quality Plan}
% =============================================================================

\subsection{Core Data Quality Checks}

\begin{itemize}[nosep]
  \item Required fields are present;
  \item event timestamps are parseable and in UTC;
  \item match IDs are non-null;
  \item \texttt{point\_winner} is one of the valid players;
  \item set/game/point numbers are non-negative;
  \item duplicate event IDs are removed or ignored;
  \item invalid scoring states are flagged.
\end{itemize}

\subsection{DQ Failure Policy}

Records failing non-critical checks are quarantined at
\texttt{s3://tennis-big-data/raw/quarantine/}. Records failing critical checks are not published
to Kafka or serving tables. Rejected record counts are logged per pipeline run.

% =============================================================================
\section{Pipelines}
% =============================================================================

\subsection{Pipeline A: Historical Ingestion}

Loads raw source files into S3, validates against \texttt{contracts/point\_event\_schema.json},
and writes zone metadata. Triggered by Airflow S3KeySensor or CLI script \texttt{ingest.sh}.

\begin{importantbox}[--- Gap \#08 addressed: ingestion trigger]
Pipeline~A has a defined entry point. Option~1 (recommended): an Airflow
\texttt{S3KeySensor} watching \texttt{s3://tennis-big-data/raw/points/}. Option~2: a bash script
\texttt{ingest.sh} that syncs a local data folder and then triggers the Airflow DAG via the CLI
(\texttt{airflow dags trigger nightly\_batch\_pipeline}). Both options are documented in
\texttt{README.md} under Quick Start.
\end{importantbox}

\subsection{Pipeline B: Curated Dataset Build}
Reads raw S3 data, standardizes columns, normalizes types, validates records, and writes cleaned
Iceberg tables.

\subsection{Pipeline C: Feature Build}
Builds batch feature tables using temporal-join-correct rolling windows. Applies point-in-time
split described in Section~15.2.

\subsection{Pipeline D: Baseline Build}
Computes player historical baselines: serve win rate, return performance, pressure-point
conversion, rally length distribution, and surface-specific tendencies. Outputs stored as
versioned Parquet in \texttt{baselines/}.

\subsection{Pipeline E: Model Training}
Trains odds and risk models using feature tables and baseline profiles. Evaluates metrics against
thresholds. Calls \texttt{publish\_model.py} for two-phase validated publish.

\subsection{Pipeline F: Replay Preparation}
Converts historical point sequences into replay-ready event files. Assigns synthetic match IDs,
replay timestamps, and partition keys.

\subsection{Pipeline G: Live Replay and Streaming}
Starts the replay producer, writes events to Kafka, runs Spark Structured Streaming, scores
events, and writes live outputs to PostgreSQL and optional Redis.

\subsection{Pipeline H: Serving API}
Always-running FastAPI service. Queries PostgreSQL and Redis, formats responses, returns
dashboard-ready JSON.

\subsection{Pipeline I: Reporting and Evidence Export}
Exports benchmark results, example match timelines, model evaluation tables, and final report
figures into S3 and the repository's \texttt{reports/} folder.

\subsection{Airflow DAG --- Explicit Dependency Graph}

\begin{criticalbox}[--- Gap \#06 addressed: DAG dependency prevents replay before model is ready]
The single MVP Airflow DAG must encode the following task dependency chain:

\begin{lstlisting}[language=python, caption={nightly\_batch\_pipeline.py --- task ordering}]
ingest_task \
  >> curated_build_task \
  >> feature_build_task \
  >> baseline_build_task \
  >> train_models_task \
  >> publish_model_task \   # validates metrics, writes latest.json
  >> replay_prep_task \     # CANNOT start before model is published
  >> notify_streaming_ready_task
\end{lstlisting}

\texttt{publish\_model\_task} writes \texttt{latest.json} only if validation passes (Section~14.4).
The streaming startup script checks for the existence of \texttt{latest.json} before subscribing
to Kafka:

\begin{lstlisting}[language=python, caption={streaming\_processor.py --- model readiness check}]
import boto3, json, sys

s3 = boto3.client('s3')
try:
    obj = s3.get_object(Bucket='tennis-big-data',
                        Key='models/odds/latest.json')
    model_meta = json.loads(obj['Body'].read())
    MODEL_PATH = model_meta['artifact_path']
except s3.exceptions.NoSuchKey:
    print("ERROR: models/odds/latest.json not found. "
          "Run Pipeline E before starting streaming.")
    sys.exit(1)
\end{lstlisting}
\end{criticalbox}

% =============================================================================
\section{Streaming Correctness and Reliability Semantics}
% =============================================================================

\subsection{Core Streaming Assumptions}

The stream is generated from historical replay so event order is controlled. Kafka partitions use
\texttt{match\_id} as the key so point events for the same match preserve order within a
partition.

\subsection{Kafka Topic Configuration}

\begin{criticalbox}[--- Gap \#01 addressed: explicit Kafka topic specification]
Before running \texttt{kafka-topics --create}, the following parameters must be decided and
committed to \texttt{infra/ec2/kafka\_setup.sh}:

\begin{lstlisting}[language=bash, caption={kafka\_setup.sh}]
KAFKA_BIN=/opt/kafka/bin
BOOTSTRAP=localhost:9092
TOPIC=tennis-point-events

$KAFKA_BIN/kafka-topics.sh --create \
  --bootstrap-server $BOOTSTRAP \
  --topic $TOPIC \
  --partitions 20 \           # supports up to 20 concurrent matches
  --replication-factor 1 \    # single-broker EC2 demo
  --config retention.ms=86400000 \   # 24-hour retention
  --config cleanup.policy=delete \
  --config max.message.bytes=1048576

# Consumer group used by Spark streaming
# (set in streaming_processor.py via .option("kafka.group.id","spark-tennis"))
echo "Topic $TOPIC created. Consumer group: spark-tennis"
\end{lstlisting}

\textbf{Partition count rationale:} 20 partitions supports up to 20 concurrent replay matches
with strict ordering per match (one partition per match\_id). Kafka does not allow safe
partition-count changes on live topics, so this must be set correctly before the first producer
run.
\end{criticalbox}

\subsection{Spark Micro-Batch Tuning}

\begin{criticalbox}[--- Gap \#02 addressed: throughput tuning parameters]
The following parameters must be set before benchmarking. Expose all three as environment
variables so they can be tuned without redeployment:

\begin{lstlisting}[language=python, caption={streaming\_processor.py --- tuning config}]
spark = SparkSession.builder \
    .appName("TennisStreamingProcessor") \
    .config("spark.sql.shuffle.partitions", "4")  \  # = vCPU count on m5.xlarge
    .getOrCreate()

query = (scored_stream
    .writeStream
    .trigger(processingTime="2 seconds")   # micro-batch interval
    .option("maxOffsetsPerTrigger", os.getenv("MAX_OFFSETS", "5000"))
    .option("checkpointLocation", "/mnt/checkpoints/streaming_processor")
    .option("kafka.group.id", "spark-tennis")
    .foreachBatch(write_batch)
    .start())
\end{lstlisting}

\textbf{Validation:} the throughput benchmark (Section~24) must demonstrate $\geq 1{,}000$
events/second by tuning \texttt{MAX\_OFFSETS} and \texttt{processingTime} together. Start with
the values above and increase \texttt{MAX\_OFFSETS} if throughput is below target.
\end{criticalbox}

\subsection{Watermarking and Late Data}

\begin{importantbox}[--- Gap \#11 addressed: watermark interval quantified]
Because replay events are generated by the project, late data is rare. The watermark is set to
30 seconds. This bounds state growth for all stateful aggregations:

\begin{lstlisting}[language=python]
windowed = (events_df
    .withWatermark("event_ts", "30 seconds")
    .groupBy(
        F.window("event_ts", "5 minutes", "30 seconds"),
        "match_id", "player_a", "player_b"
    )
    .agg(...feature_aggregations...))
\end{lstlisting}

State TTL is bounded by the watermark plus window duration. For a 5-minute window with a
30-second watermark, state is held for at most 5.5 minutes per match. After a match ends, state
is evicted automatically. Set \texttt{spark.sql.streaming.stateStore.minDeltasForSnapshot=10}
to control state snapshot frequency.
\end{importantbox}

\subsection{Checkpoint Storage}

\begin{importantbox}[--- Gap \#12 addressed: checkpoint location decision]
Use \textbf{local EBS} for checkpointing during the demo:
\texttt{checkpointLocation=/mnt/checkpoints/streaming\_processor}.

Mount a dedicated EBS volume at \texttt{/mnt} so checkpoints survive instance reboots but are
not lost on instance termination. \textbf{Important:} when the Spark streaming schema changes
(e.g.\ new feature column added), the existing checkpoint must be deleted before restart.
Document this in \texttt{README.md}:

\begin{lstlisting}[language=bash]
# Run this ONLY when streaming schema changes
sudo rm -rf /mnt/checkpoints/streaming_processor
echo "Checkpoint cleared. Restart streaming_processor.py"
\end{lstlisting}
\end{importantbox}

\subsection{Dead-Letter Queue for Malformed Messages}

\begin{criticalbox}[--- Gap \#04 addressed: streaming dead-letter handling]
\texttt{dead\_letter.py} wraps Spark's \texttt{from\_json} so a single malformed Kafka message
cannot stall the entire streaming job:

\begin{lstlisting}[language=python, caption={dead\_letter.py}]
from pyspark.sql import functions as F
from pyspark.sql.types import StructType

def parse_with_dead_letter(raw_df, schema: StructType,
                            s3_quarantine: str) -> tuple:
    """
    Returns (good_df, bad_df).
    good_df: rows where JSON parsed successfully.
    bad_df:  rows where parse failed (written to S3 quarantine).
    """
    parsed = raw_df.select(
        F.col("offset"),
        F.col("timestamp").alias("kafka_ts"),
        F.from_json(F.col("value").cast("string"),
                    schema,
                    {"mode": "PERMISSIVE"}).alias("data"),
        F.col("value").cast("string").alias("raw_value")
    )

    good_df = parsed.filter(F.col("data.event_id").isNotNull()) \
                    .select("data.*", "kafka_ts")

    bad_df  = parsed.filter(F.col("data.event_id").isNull()) \
                    .select("raw_value", "kafka_ts",
                            F.current_timestamp().alias("quarantine_ts"))

    bad_df.write \
          .mode("append") \
          .json(s3_quarantine)   # s3://tennis-big-data/raw/quarantine/streaming/

    return good_df, bad_df
\end{lstlisting}

Call \texttt{parse\_with\_dead\_letter} at the start of every micro-batch. Log the count of
quarantined records per batch to the observability metrics (Section~22).
\end{criticalbox}

\subsection{Idempotency and Deduplication}

Every event includes a unique \texttt{event\_id}. The serving writer uses PostgreSQL upserts for
\texttt{current\_match\_state} and unique constraints on \texttt{risk\_events.risk\_event\_id}.

\subsection{State TTL and Cleanup}

Streaming state evicted after watermark window expires. Redis keys carry explicit TTLs
(Section~12.2). Old demo data archived to S3 by Pipeline~I.

\subsection{Failure Handling}

\begin{itemize}[nosep]
  \item Producer retry logic on Kafka send failure;
  \item Kafka topic durability with 24-hour retention;
  \item Spark checkpoint recovery on EBS;
  \item database upserts for idempotent writes;
  \item Airflow task retries for batch jobs;
  \item dead-letter queue for malformed streaming records;
  \item logs for all failed records and failed tasks.
\end{itemize}

% =============================================================================
\section{Backend Service Plan}
% =============================================================================

\subsection{Replay Producer}

Reads prepared point-event files and publishes to Kafka. Controls replay speed, number of
concurrent matches, and timestamp generation. Uses \texttt{match\_id} as the Kafka message key
to enforce partition routing.

\subsection{Streaming Processor}

Consumes Kafka events, calls \texttt{parse\_with\_dead\_letter}, computes rolling features,
applies the validated model, and writes to PostgreSQL and optional Redis.

\subsection{API Service}

FastAPI service exposing dashboard endpoints. No heavy analytics performed inline.

\begin{importantbox}[--- Gap \#14 addressed: API bearer token authentication]
All API endpoints require a static bearer token validated in
\texttt{api/app/middleware/auth.py}:

\begin{lstlisting}[language=python, caption={auth.py}]
import os
from fastapi import Request, HTTPException
from starlette.middleware.base import BaseHTTPMiddleware

class BearerTokenMiddleware(BaseHTTPMiddleware):
    async def dispatch(self, request: Request, call_next):
        if request.url.path == "/health":
            return await call_next(request)   # health check is public
        token = request.headers.get("X-API-Key", "")
        if token != os.environ["API_KEY"]:
            raise HTTPException(status_code=401,
                                detail="Invalid or missing API key")
        return await call_next(request)
\end{lstlisting}

The \texttt{API\_KEY} value is stored in AWS SSM Parameter Store and injected at startup. The
React dashboard sends the key in every request header. Document the demo key in
\texttt{.env.example} as \texttt{API\_KEY=\$\{SSM:/tennis/api\_key\}}.
\end{importantbox}

% =============================================================================
\section{Frontend Build Plan}
% =============================================================================

\subsection{Pages}

React application pages: Overview Dashboard, Live Odds Board, Match Drill-Down, Risk Explorer,
System Benchmark View.

\subsection{Overview Dashboard Components}

Number of live matches, events processed per second, recent risk alerts, latest model version,
system health cards.

\subsection{Live Odds Board Components}

Active match table, player win probabilities, current score, momentum indicator, last update
timestamp.

\subsection{Match Drill-Down Components}

Odds history line chart, recent point timeline, rolling serve and return statistics, risk event
markers, model explanation summary.

\subsection{Risk Explorer Components}

Risk event table, filters by player/match/risk level/time, risk score distribution, feature
deviation details, explanation panel.

\subsection{System Benchmark View}

Throughput chart, latency summary, batch runtime, processed event count, resource utilization.

\subsection{Frontend Refresh Policy}

Live endpoints polled every 2--5 seconds. Historical pages refresh on user action only.

\begin{minorbox}[--- Gap \#22 addressed: Cache-Control headers]
Adding \texttt{Cache-Control} headers to FastAPI responses reduces redundant PostgreSQL reads
when multiple browser tabs or demo viewers are connected simultaneously:

\begin{lstlisting}[language=python, caption={routes/live.py --- cache headers}]
from fastapi import Response

@router.get("/api/matches/live")
async def get_live_matches(response: Response):
    response.headers["Cache-Control"] = "max-age=2, must-revalidate"
    return await db.fetch_live_matches()

@router.get("/api/matches/{match_id}/odds-history")
async def get_odds_history(match_id: str, response: Response):
    response.headers["Cache-Control"] = "max-age=30"
    return await db.fetch_odds_history(match_id)
\end{lstlisting}
\end{minorbox}

% =============================================================================
\section{API Contracts}
% =============================================================================

\subsection{Required Endpoints}

\begin{lstlisting}
GET  /health                                  (public, no auth)
GET  /api/matches/live                        Cache-Control: max-age=2
GET  /api/matches/{match_id}                  Cache-Control: max-age=5
GET  /api/matches/{match_id}/odds-history     Cache-Control: max-age=30
GET  /api/matches/{match_id}/points           Cache-Control: max-age=5
GET  /api/risk-events                         Cache-Control: max-age=5
GET  /api/risk-events/{risk_event_id}         Cache-Control: max-age=30
GET  /api/benchmarks/latest                   Cache-Control: max-age=60
GET  /api/models/current                      Cache-Control: max-age=60
\end{lstlisting}

\subsection{API Correctness Rules}

\begin{itemize}[nosep]
  \item API responses must include timestamps.
  \item Empty states return valid JSON, not errors.
  \item Database errors are handled gracefully with 500 status and logged.
  \item Dashboard endpoints optimized for repeated reads (indexed queries only).
  \item No endpoint triggers heavy Spark processing directly.
  \item All endpoints except \texttt{/health} require \texttt{X-API-Key} header.
\end{itemize}

% =============================================================================
\section{Observability Plan}
% =============================================================================

\subsection{Metrics to Capture}

\begin{itemize}[nosep]
  \item Kafka messages produced per second;
  \item Kafka messages consumed per second;
  \item Spark micro-batch duration;
  \item dead-letter message count per micro-batch;
  \item end-to-end event latency;
  \item number of active matches;
  \item number of risk alerts generated;
  \item API request latency;
  \item PostgreSQL row counts per table;
  \item batch job runtime per pipeline.
\end{itemize}

\subsection{Logs to Capture}

Producer logs, Kafka container logs, Spark driver and executor logs, Airflow task logs, FastAPI
request/error logs, frontend deployment logs, database migration logs, dead-letter quarantine
counts.

\subsection{Operational Dashboards}

The frontend Benchmark View serves as the main operational dashboard. Container-level metrics may
be captured via \texttt{docker stats} or CloudWatch if time allows.

% =============================================================================
\section{Testing Strategy}
% =============================================================================

\subsection{Unit Tests}

\begin{itemize}[nosep]
  \item Data contract validation (each schema field);
  \item feature calculations (rolling window correctness, null handling);
  \item odds scoring functions;
  \item risk scoring functions (z-score thresholds);
  \item \texttt{parse\_with\_dead\_letter} with intentionally malformed JSON;
  \item two-phase model publish --- mock S3 put, assert latest.json updated only on pass;
  \item API response formatting;
  \item bearer token middleware (valid, missing, wrong token);
  \item frontend utility functions.
\end{itemize}

\subsection{Integration Tests}

\begin{itemize}[nosep]
  \item Producer to Kafka;
  \item Spark reading from Kafka;
  \item Spark writing to PostgreSQL (including JDBC batch size);
  \item API reading from PostgreSQL;
  \item frontend calling API endpoints;
  \item dead-letter route: malformed message reaches S3 quarantine.
\end{itemize}

\subsection{End-to-End Tests}

Run a small replay from S3 sample data through Kafka, Spark, PostgreSQL, API, and dashboard.
Verify that a live match appears in the frontend and odds values update. Verify that a
deliberately malformed event does not crash the streaming job.

\subsection{Performance Tests}

Vary replay speed, number of concurrent matches, and event volume. Outputs include throughput,
latency, and resource utilization. \textbf{Discard the first 60 seconds of metrics} as the
Spark JIT, Kafka consumer rebalance, and JVM GC warm-up window (see Section~24).

% =============================================================================
\section{Benchmarking Plan}
% =============================================================================

\subsection{Benchmark Scenarios}

\begin{table}[H]
\centering
\begin{tabularx}{\linewidth}{L{4cm}X}
\toprule
\textbf{Scenario} & \textbf{Description} \\
\midrule
Baseline replay & Small number of matches, low replay rate \\
Moderate load & Hundreds of concurrent simulated matches \\
High load & Maximum sustainable replay rate for selected EC2 instance \\
Failure recovery & Restart streaming job, verify checkpoint recovery \\
Serving stress & Repeated dashboard/API reads during replay \\
\bottomrule
\end{tabularx}
\caption{Benchmark scenarios}
\end{table}

\subsection{Metrics Reported}

Events produced/second, events processed/second, average and p95 streaming latency, API response
latency, batch job runtime, model metrics, alerts generated.

\subsection{Benchmark Success Criteria}

\begin{itemize}[nosep]
  \item System processes a live replay stream end-to-end.
  \item Dashboard updates within 5 seconds of event production.
  \item Risk events stored and displayed correctly.
  \item Benchmark evidence exported to \texttt{results/benchmarks/}.
\end{itemize}

\begin{importantbox}[--- Gap \#21 addressed: warm-up period]
All benchmark measurement windows must \textbf{exclude the first 60 seconds} after streaming
job startup. This accounts for Spark JIT compilation, Kafka consumer group rebalancing, and JVM
GC ramp-up. Document the warm-up window in \texttt{results/benchmarks/README.md} so graders
understand the methodology. The benchmark script must record:
\begin{itemize}[nosep]
  \item \texttt{warmup\_seconds}: 60
  \item \texttt{measurement\_start}: timestamp after warm-up
  \item \texttt{measurement\_duration}: configurable (recommend 300 seconds)
\end{itemize}
\end{importantbox}

% =============================================================================
\section{Analytical Output Plan}
% =============================================================================

\subsection{Required Analytical Questions}

\begin{itemize}[nosep]
  \item How do win probabilities change during a simulated match?
  \item Which features contribute most to odds movement?
  \item Which players or matches generate high risk scores?
  \item How does system latency change as event throughput increases?
  \item How do batch and streaming layers interact?
\end{itemize}

\subsection{Required Charts and Tables}

Live odds over time, momentum timeline, risk score distribution, feature deviation table,
throughput and latency chart, model evaluation table, architecture diagram.

\subsection{Business Interpretation Requirement}

Every analytical output must include a plain-English interpretation. A high risk score must be
explained as a statistical anomaly relative to baseline, not as proof of wrongdoing.

% =============================================================================
\section{Report, Presentation, and Deliverables Plan}
% =============================================================================

\subsection{Required Final Deliverables}

GitHub repository with code and documentation; Docker Compose files for both EC2 instances and
local dev; S3 data lake structure with zone manifests; working Kafka replay with explicit topic
config; Spark batch and streaming jobs with dead-letter handling; trained model artifacts with
two-phase publish; PostgreSQL serving schema with indexes and archival; FastAPI backend with auth;
React frontend; benchmark results with warm-up methodology documented; final report; final
presentation.

\subsection{Project Report Structure}

Problem statement; architecture; data engineering design; Kafka and streaming configuration;
modeling methodology (unsupervised anomaly approach); implementation details; benchmark results;
dashboard screenshots; limitations and future work.

\subsection{Presentation Slide Structure}

\begin{enumerate}[nosep]
  \item Title and motivation
  \item Problem statement
  \item Architecture overview
  \item Data lake and pipelines
  \item Streaming demo flow
  \item Model design
  \item Dashboard screenshots
  \item Benchmark results
  \item Limitations
  \item Conclusion
\end{enumerate}

\subsection{Executive Audience Standard}

The final presentation explains the system clearly to both technical and business audiences. It
emphasizes: what was built, why the architecture is appropriate, what tradeoffs were made, and
what evidence proves the system works.

% =============================================================================
\section{Milestones}
% =============================================================================

\subsection{Milestone 1: Data Foundation}
Upload raw datasets to S3, define data contracts with JSON schemas, implement cleaning scripts
with schema validation, write zone manifests.

\subsection{Milestone 2: Offline Modeling}
Build features with temporal correctness, compute baselines, train initial odds and risk models,
validate against thresholds, publish via two-phase publish.

\subsection{Milestone 3: Live Replay}
Create Kafka topic with explicit config, implement producer, verify point events published and
consumed, verify dead-letter handling.

\subsection{Milestone 4: Serving Layer}
Create PostgreSQL schema with indexes and archival, implement writers with JDBC batch size,
validate current match state updates.

\subsection{Milestone 5: Frontend MVP}
Build live match overview and odds board with cache headers.

\subsection{Milestone 6: Risk Explorer}
Add risk-event API and frontend risk explorer page.

\subsection{Milestone 7: Benchmark and Report Assets}
Run performance tests with warm-up excluded, export charts and tables.

\subsection{Milestone 8: Final Demo Readiness}
Run full end-to-end demo, clean repository, finalize report and presentation.

% =============================================================================
\section{MVP vs.\ Stretch Goals}
% =============================================================================

\subsection{MVP}

S3 data lake with Iceberg on curated/features zones (Hadoop catalog); one Kafka topic with
explicit partition config; one Spark streaming job with dead-letter handling; basic odds model
with validation gate; basic risk score (z-score baseline); PostgreSQL serving tables with indexes
and archival; FastAPI endpoints with auth; React dashboard with cache headers; benchmark evidence
with warm-up documented.

\subsection{Stretch Goals}

Redis caching; Airflow multi-DAG orchestration; advanced model comparison; model version
dashboard; larger synthetic load tests; richer match drill-down visualizations; automatic report
exports; CloudWatch metrics integration.

% =============================================================================
\section{Explicit Coding Order}
% =============================================================================

\begin{enumerate}[nosep]
  \item Define data contracts and write JSON schema files
  \item Create S3 folder layout with \texttt{create\_layout.sh} and zone manifests
  \item Write IAM role policies for EC2-1 and EC2-2
  \item Configure SSM Parameter Store with all secrets
  \item Build local sample data loader with schema validation
  \item Build cleaning pipeline with Iceberg Hadoop catalog config
  \item Build feature pipeline with temporal correctness
  \item Create PostgreSQL schema with indexes and archival job
  \item Configure Kafka topic with \texttt{kafka\_setup.sh}
  \item Build replay producer
  \item Build Spark streaming read from Kafka with dead-letter handling
  \item Add rolling feature computation with watermark
  \item Add odds and risk scoring with model readiness check
  \item Add two-phase model publish with validation gate
  \item Encode Airflow DAG dependencies (E $\to$ F)
  \item Write outputs to PostgreSQL with JDBC batch size
  \item Build FastAPI endpoints with auth middleware and cache headers
  \item Build React dashboard
  \item Run benchmarks with warm-up exclusion
  \item Finalize report and demo
\end{enumerate}

% =============================================================================
\section{Code Quality, Reproducibility, and CI}
% =============================================================================

\subsection{Quality Controls}

\begin{itemize}[nosep]
  \item Consistent Python formatting (ruff or black);
  \item SSM-backed environment variable management;
  \item schema validation on all raw ingestion;
  \item unit tests for all core logic including dead-letter and auth;
  \item clear README with Quick Start, schema change instructions, checkpoint clear procedure;
  \item reproducible Docker Compose commands.
\end{itemize}

\subsection{Reproducibility Rules}

\begin{itemize}[nosep]
  \item Raw data must not be manually edited;
  \item every pipeline output reproducible from source inputs;
  \item model artifacts include version, training timestamp, and metrics;
  \item benchmark runs record configuration values including warm-up window.
\end{itemize}

\subsection{CI Pipeline}

\begin{importantbox}[--- Gap \#17 addressed: concrete CI toolchain]
\textbf{GitHub Actions} (\texttt{.github/workflows/ci.yml}) with three jobs:

\begin{lstlisting}[ caption={ci.yml}]
name: CI

on: [push, pull_request]

jobs:
  lint:
    runs-on: ubuntu-latest
    steps:
      - uses: actions/checkout@v4
      - uses: actions/setup-python@v5
        with: {python-version: "3.11"}
      - run: pip install ruff
      - run: ruff check spark/ api/ producer/

  test:
    runs-on: ubuntu-latest
    steps:
      - uses: actions/checkout@v4
      - uses: actions/setup-python@v5
        with: {python-version: "3.11"}
      - run: pip install -r requirements-test.txt
      - run: pytest tests/unit/ --tb=short -q

  compose-validate:
    runs-on: ubuntu-latest
    steps:
      - uses: actions/checkout@v4
      - run: docker compose -f docker-compose.processing.yml config
      - run: docker compose -f docker-compose.serving.yml config
      - run: docker compose -f docker-compose.dev.yml config
\end{lstlisting}

This runs in under 3 minutes and catches formatting failures, broken unit tests, and invalid
Docker Compose configs before they reach the EC2 instances.
\end{importantbox}

% =============================================================================
\section{Environment Variables}
% =============================================================================

All secrets stored in AWS SSM Parameter Store. Non-secret configuration stored in \texttt{.env}
files on each EC2 instance, bootstrapped from \texttt{.env.example}.

\begin{lstlisting}[caption={.env.example}]
# AWS
AWS_REGION=us-east-1
S3_BUCKET=tennis-big-data

# Kafka (EC2-1 internal)
KAFKA_BOOTSTRAP_SERVERS=localhost:9092
KAFKA_TOPIC=tennis-point-events
KAFKA_PARTITIONS=20
KAFKA_CONSUMER_GROUP=spark-tennis

# Spark tuning
SPARK_MAX_OFFSETS=5000
SPARK_TRIGGER_INTERVAL=2 seconds
SPARK_SHUFFLE_PARTITIONS=4

# PostgreSQL (read from SSM on EC2-2, referenced by private IP on EC2-1)
POSTGRES_HOST=${SSM:/tennis/postgres_host}
POSTGRES_PORT=5432
POSTGRES_DB=tennis_serving
POSTGRES_USER=tennis_user
POSTGRES_PASSWORD=${SSM:/tennis/postgres_password}

# Redis
REDIS_HOST=${SSM:/tennis/redis_host}
REDIS_PORT=6379

# Models
MODEL_ODDS_PATH=s3://tennis-big-data/models/odds/latest.json
MODEL_RISK_PATH=s3://tennis-big-data/models/risk/latest.json

# API
API_BASE_URL=http://ec2-serving-elastic-ip:8000
API_KEY=${SSM:/tennis/api_key}

# Benchmark
BENCHMARK_WARMUP_SECONDS=60
BENCHMARK_MEASUREMENT_SECONDS=300
\end{lstlisting}

% =============================================================================
\section{Milestone Gates}
% =============================================================================

\subsection{Gate 1}
Data can be read from S3, schema-validated, and cleaned into Iceberg-managed curated Parquet.

\subsection{Gate 2}
Offline features (with temporal correctness) and model artifacts (passing validation thresholds)
are generated and published via two-phase publish.

\subsection{Gate 3}
Kafka topic created with explicit config; replay and Spark streaming produce scored outputs;
dead-letter handling verified with malformed test message.

\subsection{Gate 4}
Dashboard displays live odds and risk events from serving tables; archival job confirmed running;
API auth verified.

% =============================================================================
\section{Course Presentation Strategy}
% =============================================================================

The presentation focuses on architecture first, then demonstrates evidence. The strongest story
is that one event stream powers two analytical products: live odds and integrity monitoring. The
system is explained as a practical Lambda-style architecture where S3 stores durable data, EC2-1
performs batch and streaming computation, and EC2-2 serves low-latency application views.

Emphasize the gap-resolution decisions as evidence of engineering rigor: explicit Kafka topic
sizing, validated model publish, dead-letter handling, IAM separation, and temporal feature
correctness.

% =============================================================================
\section{Limitations}
% =============================================================================

\begin{itemize}[nosep]
  \item The system uses simulated live streams, not actual live tennis feeds.
  \item Risk scores are statistical indicators, not evidence of misconduct.
  \item The two-EC2 deployment is not production scale.
  \item Model accuracy depends on available historical data quality.
  \item Synthetic amplification is useful for systems testing but does not create new
    real-world evidence.
  \item The Hadoop Iceberg catalog is not suitable for concurrent multi-writer production
    workloads; for production, AWS Glue catalog would replace it.
  \item The static bearer token API auth is sufficient for demo purposes only; production
    would use OAuth~2.0 or API gateway-level auth.
\end{itemize}

% =============================================================================
\appendix
\section*{Appendix A: AI Coding Agent Execution Notes}
\addcontentsline{toc}{section}{Appendix A: AI Coding Agent Execution Notes}
% =============================================================================

When using an AI coding assistant, request small, testable components rather than the full system
at once. Each generated module must be validated with unit tests or sample input/output checks.
The AI assistant must not change architecture decisions unless the team explicitly approves.

Recommended sequence: data contracts and schemas; cleaning functions; producer logic; Spark
transformations (batch then streaming with dead-letter); API endpoints with auth; frontend
components; deployment scripts.

% =============================================================================
\section*{Appendix B: Architecture and Flow Diagrams}
\addcontentsline{toc}{section}{Appendix B: Architecture and Flow Diagrams}
% =============================================================================

\subsection*{B.1 Corrected High-Level System Architecture}

\begin{center}
\begin{tabular}{|c|}
\hline
\textbf{Amazon S3 Data Lake}\\
raw / cleaned (Iceberg) / features (Iceberg) / baselines / models / stream\_archive / results\\
\hline
\end{tabular}

\vspace{6pt}
$\updownarrow$ features, models, outputs \hspace{2cm} $\updownarrow$ scored results

\vspace{6pt}
\begin{tabular}{|c|c|c|}
\hline
\textbf{EC2-1 Processing Node} & $\longrightarrow$ & \textbf{EC2-2 Serving/App Node}\\
Kafka + Spark + Airflow + ML & scored outputs & PostgreSQL + Redis + API + React\\
\hline
\end{tabular}

\vspace{6pt}
$\uparrow$ Replay Producer \hspace{4cm} $\downarrow$ Dashboard users (Live Odds + Risk Explorer)
\end{center}

\subsection*{B.2 Corrected Data Lifecycle and Storage Flow}

\begin{enumerate}[nosep]
  \item Raw data uploaded to S3 raw zone (schema-validated).
  \item Cleaning jobs produce curated Iceberg tables.
  \item Feature jobs produce feature Iceberg tables (temporally correct).
  \item Baseline jobs produce player profiles.
  \item Model training jobs publish artifacts via two-phase validated publish.
  \item Replay jobs create live-like events.
  \item Streaming outputs archived back to S3 stream\_archive zone.
\end{enumerate}

\subsection*{B.3 Offline Training and Model Publication Flow}

Cleaned Data (S3 Iceberg) $\to$ Feature Build (Spark Batch) $\to$ Baseline Build $\to$
Model Training $\to$ Validation Gate (AUC $\geq 0.55$) $\to$ Publish to S3 + update
\texttt{latest.json} $\to$ Airflow notifies replay prep.

\subsection*{B.4 Real-Time Streaming and Scoring Flow}

Producer $\to$ Kafka (\texttt{tennis-point-events}, 20 partitions) $\to$
Spark Structured Streaming (dead-letter split) $\to$ rolling features (watermark 30s) $\to$
model load from \texttt{latest.json} $\to$ PostgreSQL upserts (JDBC batch 500) $\to$
optional Redis TTL cache.

\subsection*{B.5 MVP vs.\ Stretch Architecture Scope}

\begin{table}[H]
\centering
\begin{tabularx}{\linewidth}{L{4cm}L{5cm}X}
\toprule
\textbf{Area} & \textbf{MVP} & \textbf{Stretch} \\
\midrule
Data Lake & S3 zones + Iceberg (Hadoop catalog) + zone manifests & Glue catalog, richer partitioning \\
Streaming & One Kafka topic, 20 partitions, dead-letter & Multiple topics, retry queues \\
Processing & One Spark streaming job & Separate odds and risk streams \\
Models & Baseline odds/risk with validation gate & Advanced model comparison \\
Serving & PostgreSQL + indexes + archival & PostgreSQL + optimized Redis cache \\
Frontend & Dashboard MVP + cache headers & Rich drill-down and explainability \\
Orchestration & Single Airflow DAG with explicit dependencies & Multi-DAG pipeline suite \\
Auth & Static bearer token & OAuth 2.0 \\
Secrets & AWS SSM Parameter Store & HashiCorp Vault \\
\bottomrule
\end{tabularx}
\caption{MVP vs.\ stretch architecture comparison}
\end{table}

\subsection*{B.6 Gap Resolution Summary}

\begin{longtable}{L{1cm}L{3.5cm}L{2.5cm}L{7cm}}
\toprule
\textbf{Gap} & \textbf{Title} & \textbf{Severity} & \textbf{Section Addressed} \\
\midrule
\endhead
\#01 & Kafka topic config & Critical & Section~18.2 \\
\#02 & Spark micro-batch tuning & Critical & Section~18.2 \\
\#03 & PostgreSQL write bottleneck & Critical & Section~12.1 \\
\#04 & Dead-letter queue & Critical & Section~18.5 \\
\#05 & IAM role separation & Critical & Section~8.3 \\
\#06 & Airflow DAG dependency & Critical & Section~17.10 \\
\#07 & Point-in-time feature correctness & Critical & Section~15.2 \\
\#08 & Ingestion trigger & Important & Section~17.1 \\
\#09 & Iceberg catalog backend & Important & Section~7.2 \\
\#10 & EC2-1 OOM risk & Important & Section~8.2 \\
\#11 & Watermark interval & Important & Section~18.3 \\
\#12 & Checkpoint location & Important & Section~18.4 \\
\#13 & Two-phase model publish & Important & Section~14.4 \\
\#14 & FastAPI auth & Important & Section~19.3 \\
\#15 & odds\_history size & Important & Section~12.3 \\
\#16 & Anomaly training strategy & Important & Section~13.2 \\
\#17 & CI toolchain & Important & Section~30.3 \\
\#18 & Shared secrets management & Important & Section~8.4 \\
\#19 & S3 zone manifests & Minor & Section~10.2 \\
\#20 & Model evaluation thresholds & Minor & Section~14.4 \\
\#21 & Benchmark warm-up period & Minor & Section~24.3 \\
\#22 & API cache headers & Minor & Section~20.7 \\
\#23 & Docker Compose network & Minor & Section~8.1 \\
\#24 & Parquet schema evolution & Minor & Section~9.4 \\
\bottomrule
\caption{Complete gap resolution index}
\end{longtable}

\end{document}
